% Source: https://artofproblemsolving.com/community/c3h3237635_50k_message_milestone
% Stable identity: aops-topic-3237635-q09
The point $X$ is the incenter. Put $A=(-12,0),C=(12,0)$ and $B=(0,6\sqrt5)$, since the altitude is $\sqrt{18^2-12^2}=6\sqrt5$. The triangle area is $72\sqrt5$ and semiperimeter $30$, so its inradius is $r=\dfrac{12\sqrt5}{5}$. By symmetry $X=(0,r)$. Therefore
\[BX=6\sqrt5-r=\dfrac{18\sqrt5}{5},\qquad CX=\sqrt{12^2+r^2}=\dfrac{12\sqrt{30}}{5}.\]
Adding $BC=18$ gives $\boxed{18+\dfrac{18\sqrt5+12\sqrt{30}}{5}}$.
\begin{asy}
usepackage("amsmath");
size(7cm);
pair A=(-12,0),C=(12,0),B=(0,6*sqrt(5)),X=(0,12*sqrt(5)/5);draw(A--B--C--cycle);draw(A--X--C,dashed);draw(B--X,blue);draw(circle(X,12*sqrt(5)/5),gray);label("$A$",A,SW);label("$B$",B,N);label("$C$",C,SE);label("$X$",X,E);label("$24$",(0,0),S);label("$18$",(A+B)/2,NW);label("$18$",(C+B)/2,NE);
\end{asy}
